\clearpage\begingroup\SourceLeaf{41}\HistoricalHead{36}{L}{PHYSIK. BRIEFWECHSEL.}
\begin{center}
[3.]\\[.7em]
\textsc{Gerling} an \textsc{Gauss}. Marburg, 15. Februar 1852.
\end{center}
\noindent .\ .\ .\ .\ .\ .\ .\ .\ .\ .\ .\ .\ .\ .\ .\ .\ .\ .\ .\ .\ .
\par\noindent\hspace*{\parindent}\{Wegen des Pendels schreibe ich ein andermal, so Gott will, weitläuftiger.
Vor der Hand will ich nur die \textls[120]{sichere Thatsache} bemerken, dass ein \textls[120]{gehörig}
eingeklemmter Metalldraht die Kugel durch seine Torsion immer so
dreht, dass ihre eigenen Verticalkreise immer in ihren ursprünglichen Azimuthen
bleiben. Ich habe dieses in allen Azimuthen der Schwingungs-Ebene,
die sich mitunter 20, ja 30 Grad während der Versuche änderte, bewährt
gefunden.\}\quad .\ .\ .\ .\ .\ .\ .\ .\ .\ .\ .\ .\ .\ .\ .\ .\ .\ .\ .\ .\ .
\clearpage\endgroup
